% GENERATED FILE. Edit canonical lessons, then run npm run generate:books.
% canonical-source-hash: fnv1a64:ae47ef94c3668f74
% canonical-lessons: HI-C327-purvaj, HI-C327-basna, HI-C327-vams, HI-C327-yug, HI-C327-ujarna, HI-R327-review-320-327

\chapter{Ancestors}
\label{ch:hi-ancestors}
\hlchaptermodality{327}

\begin{chapteropening}
\textbf{By the end of this chapter:} \emph{I can tell the story of a family line.}

Complete the last lesson of chapter 327: I can tell the story of a family line.
\end{chapteropening}

\section[\texorpdfstring{\hi{पूर्वज} (pūrvaj)}{पूर्वज (pūrvaj)}]{\hi{पूर्वज} (pūrvaj) — an ancestor}

\label{lesson:HI-C327-purvaj}



\begin{quote}
Before the new one: say the Hindi for an inscription (on stone), then the Hindi for remains.
\end{quote}

\subsection*{You'll want to know: \hi{पूर्वज}}
\textbf{\hi{पूर्वज}} — \emph{pūrvaj} — \textquotedblleft{}an ancestor\textquotedblright{}.

\begin{quote}
\hi{हमारे} \hi{पूर्वज} \hi{सौ} \hi{साल} \hi{पहले} \hi{इस} \hi{गाँव} \hi{में} \hi{आए।}
\end{quote}

\begin{quote}
\emph{hamāre pūrvaj sau sāl pahle is gā̃v meṁ āe.}
\end{quote}

\begin{quote}
Our ancestors came to this village a hundred years ago.
\end{quote}

Masculine.

\begin{cousinweb}
\textbf{\hi{पूर्वज}} (\emph{pūrvaj}) is Sanskrit \emph{pūrvaja}, \textquotedblleft{}born before\textquotedblright{} (Monier-Williams), \emph{pūrva}, earlier, and \emph{-ja}, born: those who came before us.
\end{cousinweb}

\begin{grammarlens}[title={What you've built}]
One more word for telling what happened.
\end{grammarlens}

\subsection*{Your turn}
\begin{itemize}
\raggedright
  \item \emph{Say it:} \emph{pūrvaj}
  \item \emph{Say it:} \emph{pūrvaj}, once more
  \item \emph{Say it:} \emph{hamāre pūrvaj sau sāl pahle is gā̃v meṁ āe.} — Our ancestors came to this village a hundred years ago.
  \item \emph{Recall:} say the Hindi for an inscription (on stone), then the Hindi for remains, then say \emph{pūrvaj} again
\end{itemize}

\begin{tcolorbox}[breakable,colback=teal!4,colframe=teal!35!black,title={Before you move on}]
What does \hi{पूर्वज} mean? (\textquotedblleft{}An ancestor\textquotedblright{}.) Say it once more.
\end{tcolorbox}
\section[\texorpdfstring{\hi{बसना} (basnā)}{बसना (basnā)}]{\hi{बसना} (basnā) — to settle; to make a home}

\label{lesson:HI-C327-basna}



\begin{quote}
Before the new one: say the Hindi for remains, then the Hindi for an ancestor.
\end{quote}

\subsection*{You'll want to know: \hi{बसना}}
\textbf{\hi{बसना}} — \emph{basnā} — \textquotedblleft{}to settle; to make a home\textquotedblright{}.

\begin{quote}
\hi{वे} \hi{नदी} \hi{के} \hi{किनारे} \hi{बसे।}
\end{quote}

\begin{quote}
\emph{ve nadī ke kināre base.}
\end{quote}

\begin{quote}
They settled on the riverbank.
\end{quote}

Past (it …): \hi{वह} \hi{परिवार} \hi{यहाँ} \hi{बसा}, \emph{vah parivār yahā̃ basā}.

\begin{cousinweb}
\textbf{\hi{बसना}} (\emph{basnā}) goes back to Sanskrit \emph{vasati}, \textquotedblleft{}stays, dwells\textquotedblright{} (CDIAL 11435). A family \textbf{\hi{बसता} \hi{है}} in a place and stays; a town \textbf{\hi{बसता} \hi{है}} when people come to live there.
\end{cousinweb}

\begin{grammarlens}[title={What you've built}]
One more word for telling what happened.
\end{grammarlens}

\subsection*{Your turn}
\begin{itemize}
\raggedright
  \item \emph{Say it:} \emph{basnā}
  \item \emph{Say it:} \emph{basnā}, once more
  \item \emph{Say it:} \emph{ve nadī ke kināre base.} — They settled on the riverbank.
  \item \emph{Recall:} say the Hindi for remains, then the Hindi for an ancestor, then say \emph{basnā} again
\end{itemize}

\begin{tcolorbox}[breakable,colback=teal!4,colframe=teal!35!black,title={Before you move on}]
What does \hi{बसना} mean? (\textquotedblleft{}To settle; to make a home\textquotedblright{}.) Say it once more.
\end{tcolorbox}
\section[\texorpdfstring{\hi{वंश} (vaṁś)}{वंश (vaṁś)}]{\hi{वंश} (vaṁś) — a family line; a dynasty}

\label{lesson:HI-C327-vams}



\begin{quote}
Before the new one: say the Hindi for an ancestor, then the Hindi for to settle.
\end{quote}

\subsection*{You'll want to know: \hi{वंश}}
\textbf{\hi{वंश}} — \emph{vaṁś} — \textquotedblleft{}a family line; a dynasty\textquotedblright{}.

\begin{quote}
\hi{उस} \hi{राजा} \hi{का} \hi{वंश} \hi{तीन} \hi{सौ} \hi{साल} \hi{चला।}
\end{quote}

\begin{quote}
\emph{us rājā kā vaṁś tīn sau sāl calā.}
\end{quote}

\begin{quote}
That king's dynasty lasted three hundred years.
\end{quote}

Masculine.

\begin{cousinweb}
\textbf{\hi{वंश}} (\emph{vaṁś}) is Sanskrit \emph{vaṁśa} (Monier-Williams), a bamboo and then a lineage, one joint growing out of the last.
\end{cousinweb}

\begin{grammarlens}[title={What you've built}]
One more word for telling what happened.
\end{grammarlens}

\subsection*{Your turn}
\begin{itemize}
\raggedright
  \item \emph{Say it:} \emph{vaṁś}
  \item \emph{Say it:} \emph{vaṁś}, once more
  \item \emph{Say it:} \emph{us rājā kā vaṁś tīn sau sāl calā.} — That king's dynasty lasted three hundred years.
  \item \emph{Recall:} say the Hindi for an ancestor, then the Hindi for to settle, then say \emph{vaṁś} again
\end{itemize}

\begin{tcolorbox}[breakable,colback=teal!4,colframe=teal!35!black,title={Before you move on}]
What does \hi{वंश} mean? (\textquotedblleft{}A family line; a dynasty\textquotedblright{}.) Say it once more.
\end{tcolorbox}
\section[\texorpdfstring{\hi{युग} (yug)}{युग (yug)}]{\hi{युग} (yug) — an age; an era}

\label{lesson:HI-C327-yug}



\begin{quote}
Before the new one: say the Hindi for to settle, then the Hindi for a family line.
\end{quote}

\subsection*{You'll want to know: \hi{युग}}
\textbf{\hi{युग}} — \emph{yug} — \textquotedblleft{}an age; an era\textquotedblright{}.

\begin{quote}
\hi{वह} \hi{युग} \hi{अब} \hi{बीत} \hi{गया।}
\end{quote}

\begin{quote}
\emph{vah yug ab bīt gayā.}
\end{quote}

\begin{quote}
That age has passed now.
\end{quote}

Masculine.

\begin{cousinweb}
\textbf{\hi{युग}} (\emph{yug}) is Sanskrit \emph{yuga} (Monier-Williams), a yoke and then an age of the world. Hindi uses it for any long period: \textbf{\hi{आधुनिक} \hi{युग}}, the modern age.
\end{cousinweb}

\begin{grammarlens}[title={What you've built}]
One more word for telling what happened.
\end{grammarlens}

\subsection*{Your turn}
\begin{itemize}
\raggedright
  \item \emph{Say it:} \emph{yug}
  \item \emph{Say it:} \emph{yug}, once more
  \item \emph{Say it:} \emph{vah yug ab bīt gayā.} — That age has passed now.
  \item \emph{Recall:} say the Hindi for to settle, then the Hindi for a family line, then say \emph{yug} again
\end{itemize}

\begin{tcolorbox}[breakable,colback=teal!4,colframe=teal!35!black,title={Before you move on}]
What does \hi{युग} mean? (\textquotedblleft{}An age; an era\textquotedblright{}.) Say it once more.
\end{tcolorbox}
\section[\texorpdfstring{\hi{उजड़ना} (ujaṛnā)}{उजड़ना (ujaṛnā)}]{\hi{उजड़ना} (ujaṛnā) — to be ruined; to become deserted}

\label{lesson:HI-C327-ujarna}



\begin{quote}
Before the new one: say the Hindi for a family line, then the Hindi for an age.
\end{quote}

\subsection*{You'll want to know: \hi{उजड़ना}}
\textbf{\hi{उजड़ना}} — \emph{ujaṛnā} — \textquotedblleft{}to be ruined; to become deserted\textquotedblright{}.

\begin{quote}
\hi{अकाल} \hi{के} \hi{बाद} \hi{पूरा} \hi{गाँव} \hi{उजड़} \hi{गया।}
\end{quote}

\begin{quote}
\emph{akāl ke bād pūrā gā̃v ujaṛ gayā.}
\end{quote}

\begin{quote}
After the famine the whole village was deserted.
\end{quote}

Past (it …): \hi{वह} \hi{गाँव} \hi{उजड़} \hi{गया}, \emph{vah gā̃v ujaṛ gayā}.

\begin{cousinweb}
\textbf{\hi{उजड़ना}} (\emph{ujaṛnā}) is inherited from \emph{*ujjaṭati}, \textquotedblleft{}is uprooted\textquotedblright{}, a form CDIAL reconstructs (CDIAL 1661). A place \textbf{\hi{उजड़ता} \hi{है}} when its people leave and it falls empty.
\end{cousinweb}

\begin{grammarlens}[title={What you've built}]
One more word for telling what happened.
\end{grammarlens}

\subsection*{Your turn}
\begin{itemize}
\raggedright
  \item \emph{Say it:} \emph{ujaṛnā}
  \item \emph{Say it:} \emph{ujaṛnā}, once more
  \item \emph{Say it:} \emph{akāl ke bād pūrā gā̃v ujaṛ gayā.} — After the famine the whole village was deserted.
  \item \emph{Recall:} say the Hindi for a family line, then the Hindi for an age, then say \emph{ujaṛnā} again
\end{itemize}

\begin{tcolorbox}[breakable,colback=teal!4,colframe=teal!35!black,title={Before you move on}]
What does \hi{उजड़ना} mean? (\textquotedblleft{}To be ruined; to become deserted\textquotedblright{}.) Say it once more.
\end{tcolorbox}
\section[(dialogue)]{Looking back: from Kings and Emperors to Ancestors}

\label{lesson:HI-R327-review-320-327}



\begin{quote}
Say the Hindi for an emperor and a sovereign.
\end{quote}

\subsection*{Your turn}
Say these aloud:

\begin{itemize}
\raggedright
  \item an emperor, a sovereign, a throne, a crown, a royal court
  \item an attack, to surround, a cannon, to clash, to fight hard
  \item a sword, a shield, a spear, an arrow, for the most part
  \item to defeat, to throw down, a victory, a martyr, a treaty
  \item the subjects (of a king), an order, a slave, a rebellion, to challenge
  \item an English person, a movement (of people), a struggle, a sacrifice, to entrust
  \item ruins, to dig, a statue, an inscription (on stone), remains
  \item an ancestor, to settle, a family line, an age, to be ruined
\end{itemize}

\begin{tcolorbox}[breakable,colback=teal!4,colframe=teal!35!black,title={Before you move on}]
An emperor? (\textbf{\hi{सम्राट}}, \emph{samrāṭ}.) A sovereign? (\textbf{\hi{बादशाह}}, \emph{bādśāh}.) A throne? (\textbf{\hi{सिंहासन}}, \emph{siṁhāsan}.) A crown? (\textbf{\hi{मुकुट}}, \emph{mukuṭ}.) A royal court? (\textbf{\hi{दरबार}}, \emph{darbār}.) An attack? (\textbf{\hi{हमला}}, \emph{hamlā}.) To surround? (\textbf{\hi{घेरना}}, \emph{ghernā}.) A cannon? (\textbf{\hi{तोप}}, \emph{top}.) To clash? (\textbf{\hi{भिड़ना}}, \emph{bhiṛnā}.) To fight hard? (\textbf{\hi{जूझना}}, \emph{jūjhnā}.) A sword? (\textbf{\hi{तलवार}}, \emph{talvār}.) A shield? (\textbf{\hi{ढाल}}, \emph{ḍhāl}.) A spear? (\textbf{\hi{भाला}}, \emph{bhālā}.) An arrow? (\textbf{\hi{तीर}}, \emph{tīr}.) For the most part? (\textbf{\hi{प्रायः}}, \emph{prāyaḥ}.) To defeat? (\textbf{\hi{हराना}}, \emph{harānā}.) To throw down? (\textbf{\hi{पछाड़ना}}, \emph{pachāṛnā}.) A victory? (\textbf{\hi{विजय}}, \emph{vijay}.) A martyr? (\textbf{\hi{शहीद}}, \emph{śahīd}.) A treaty? (\textbf{\hi{संधि}}, \emph{sandhi}.) The subjects (of a king)? (\textbf{\hi{प्रजा}}, \emph{prajā}.) An order? (\textbf{\hi{हुक्म}}, \emph{hukm}.) A slave? (\textbf{\hi{ग़ुलाम}}, \emph{ġulām}.) A rebellion? (\textbf{\hi{विद्रोह}}, \emph{vidroh}.) To challenge? (\textbf{\hi{ललकारना}}, \emph{lalkārnā}.) An English person? (\textbf{\hi{अंग्रेज़}}, \emph{aṅgrez}.) A movement (of people)? (\textbf{\hi{आंदोलन}}, \emph{āndolan}.) A struggle? (\textbf{\hi{संघर्ष}}, \emph{saṅgharṣ}.) A sacrifice? (\textbf{\hi{बलिदान}}, \emph{balidān}.) To entrust? (\textbf{\hi{सौंपना}}, \emph{saumpnā}.) Ruins? (\textbf{\hi{खंडहर}}, \emph{khaṇḍhar}.) To dig? (\textbf{\hi{खोदना}}, \emph{khodnā}.) A statue? (\textbf{\hi{मूर्ति}}, \emph{mūrti}.) An inscription (on stone)? (\textbf{\hi{शिलालेख}}, \emph{śilālekh}.) Remains? (\textbf{\hi{अवशेष}}, \emph{avśeṣ}.) An ancestor? (\textbf{\hi{पूर्वज}}, \emph{pūrvaj}.) To settle? (\textbf{\hi{बसना}}, \emph{basnā}.) A family line? (\textbf{\hi{वंश}}, \emph{vaṁś}.) An age? (\textbf{\hi{युग}}, \emph{yug}.) To be ruined? (\textbf{\hi{उजड़ना}}, \emph{ujaṛnā}.)
\end{tcolorbox}
